% !TeX root = ../../Book3.tex
% 纲要标题：### 5.2 素理想的提升
% 纲要位置：工作记录/计划纲要.md，第 2042 行。
% 纲要细目（写作范围）：
% * Lying-over 定理；局部化与商域判据所确定的收缩
% * Incomparability 定理；转入分式域及同一纤维的不可比性
% * Going-up 定理；取商保证指定下端

\section{素理想的提升}\label{b3:sec:0502}

环包含 \(A\subseteq B\) 给出素谱上的收缩映射 \(\pi:\operatorname{Spec}B\to\operatorname{Spec}A\)、\(\mathfrak q\mapsto\mathfrak q\cap A\)，且收缩总保持素理想间的包含关系。整性还给出三项约束：每个底层素理想都有原像；收缩到同一个底层素理想的两个素理想若可比较，便必须相等；给定 \(\mathfrak p\subseteq\mathfrak p'\) 以及 \(\mathfrak p\) 上方的 \(\mathfrak q\)，还可找到包含 \(\mathfrak q\) 且收缩为 \(\mathfrak p'\) 的素理想。

\subsection{Lying-over 定理}\label{b3:subsec:0502:01}

在整扩张 \(A\subseteq B\) 中，若 \(\mathfrak q\) 极大，则 \(B/\mathfrak q\) 是域。要证明其收缩也极大，就要由整性推出下方的商环 \(A/(\mathfrak q\cap A)\) 也是域。下面的域判据提供了这一推理，随后将用于局部化后的扩张。

\begin{lemma}\label{b3:lem:integral-field-criterion}
设 \(D\subseteq E\) 为交换含幺整环的整扩张。则 \(D\) 是域当且仅当 \(E\) 是域。
\end{lemma}
\begin{proof}
若 \(D\) 是域，对非零 \(e\in E\) 取最低次数的首一方程。其常数项非零，否则可在整环中约去 \(e\) 而降次。把方程除以 \(e\)，再除以非零常数项，可把 \(e^{-1}\) 写成 \(e\) 的多项式，故 \(E\) 是域。

若 \(E\) 是域，对 \(0\ne d\in D\)，\(d^{-1}\in E\) 整。把它的次数 \(n\) 的首一方程乘以 \(d^{n-1}\)，得 \(d^{-1}\in D\)。所以 \(D\) 是域。
\end{proof}

设 \(A\subseteq B\) 为交换含幺环的整扩张，\(\mathfrak q\in\operatorname{Spec}B\)，\(\mathfrak p=\mathfrak q\cap A\)。收缩保证 \(\mathfrak p\) 是素理想，\cref{b3:prop:integral-quotient-localization}给出整包含
\[
A/\mathfrak p\hookrightarrow B/\mathfrak q.
\]
这两个商环都是整环，因而域判据给出
\[
\mathfrak q\;\text{极大}
\ \Longleftrightarrow\ B/\mathfrak q\;\text{是域}
\ \Longleftrightarrow\ A/\mathfrak p\;\text{是域}
\ \Longleftrightarrow\ \mathfrak p\;\text{极大}.
\]

\begin{theorem}[Lying-over]\label{b3:thm:lying-over}
设 \(A\subseteq B\) 为交换含幺环的整扩张。则对每个 \(\mathfrak p\in\operatorname{Spec}A\)，存在 \(\mathfrak q\in\operatorname{Spec}B\) 满足 \(\mathfrak q\cap A=\mathfrak p\)。这一结论称为\term[lying-over]{Lying-over 定理}[lying-over theorem]。
\end{theorem}
\begin{proof}
在 \(\mathfrak p\) 处局部化后，底环的唯一极大理想是 \(\mathfrak pA_{\mathfrak p}\)。这样，上环中任取一个极大理想，只要其收缩仍然极大，就必收缩到这个指定的理想。取 \(S=A\setminus\mathfrak p\)，\cref{b3:prop:integral-quotient-localization}给出整包含 \(A_{\mathfrak p}\subseteq S^{-1}B\)。因 \(A_{\mathfrak p}\ne0\) 且包含仍是单射，\(S^{-1}B\ne0\)，故可取其极大理想 \(\mathfrak n\)。令 \(\mathfrak m=\mathfrak n\cap A_{\mathfrak p}\)，则 \(\mathfrak m\) 为素理想；再次取商得到整包含
\[
A_{\mathfrak p}/\mathfrak m\hookrightarrow (S^{-1}B)/\mathfrak n.
\]
此时两边都是整环，右边还是域。由\cref{b3:lem:integral-field-criterion}，左边也是域，所以 \(\mathfrak m\) 极大，只能是 \(\mathfrak pA_{\mathfrak p}\)。局部化素理想对应\cref{b3:thm:spec-localization}将 \(\mathfrak n\) 收缩为 \(B\) 中不交 \(S\) 的素理想 \(\mathfrak q\)，而
\[
\begin{aligned}
\mathfrak q\cap A
&=\{\,a\in A\mid a/1\in\mathfrak n\cap A_{\mathfrak p}\,\}\\
&=\{\,a\in A\mid a/1\in\mathfrak pA_{\mathfrak p}\,\}
=\mathfrak p.
\end{aligned}
\]
\end{proof}

包含是必要条件。若只给一个不单射的整同态 \(A\to B\)，像是 \(\operatorname{Spec}A\) 中包含同态核的闭集；应对 \(A/\ker\to B\) 应用定理。

\subsection{Incomparability 定理}\label{b3:subsec:0502:02}

\begin{theorem}[Incomparability]\label{b3:thm:incomparability}
设 \(A\subseteq B\) 为交换含幺环的整扩张，\(\mathfrak q\subseteq\mathfrak q'\) 是 \(B\) 的素理想。若二者在 \(A\) 中的收缩都为 \(\mathfrak p\)，则 \(\mathfrak q=\mathfrak q'\)。这称为\term[incomparability]{Incomparability 定理}[incomparability theorem]。
\end{theorem}
\begin{proof}
先模掉 \(\mathfrak q\)，把所要证明的等式变成 \(\mathfrak q'/\mathfrak q=0\)。此时 \(A/\mathfrak p\) 嵌入整环 \(B/\mathfrak q\)；再倒置底环 \(A/\mathfrak p\) 的全部非零元素，把底环变为域 \(\operatorname{Frac}(A/\mathfrak p)\)，并得到在它上面整的整环
\[
C=(A\setminus\mathfrak p)^{-1}(B/\mathfrak q).
\]
由\cref{b3:lem:integral-field-criterion}，\(C\) 也成为域。\(\mathfrak q'/\mathfrak q\) 不交这些分母，故对应于 \(C\) 的一个素理想，而域中只有零素理想。所倒置的分母在整环 \(B/\mathfrak q\) 中都非零，局部化映射因此单射，\(\mathfrak q'/\mathfrak q\) 不可能有非零元素在局部化后消失，故它已经为零。
\end{proof}

同一底层素理想仍可有多个原像。例如整扩张 \(\mathbb Z\subseteq\mathbb Z[i]\) 中，\((2+i)\) 与 \((2-i)\) 的商环都同构于 \(\mathbb F_5\)，其中 \(i\) 分别映为 \(-2\) 与 \(2\)。两理想因此都收缩到 \((5)\)，而且都是极大理想；\(2+i\) 在第二个商环中的像为 \(4\ne0\)，所以两理想不同，便互不包含。

\subsection{Going-up 定理}\label{b3:subsec:0502:03}

\begin{theorem}[Going-up]\label{b3:thm:going-up}
设 \(A\subseteq B\) 为交换含幺环的整扩张，\(\mathfrak p\subseteq\mathfrak p'\) 是 \(A\) 的素理想，\(\mathfrak q\) 是 \(B\) 中满足 \(\mathfrak q\cap A=\mathfrak p\) 的素理想。存在 \(\mathfrak q'\supseteq\mathfrak q\) 满足 \(\mathfrak q'\cap A=\mathfrak p'\)。这称为\term[going-up]{Going-up 定理}[going-up theorem]。
\end{theorem}
\begin{proof}
要保证新素理想包含已指定的 \(\mathfrak q\)，可以把 \(\mathfrak q\) 商成零：\(B/\mathfrak q\) 中每个素理想的原像都会包含它。由\cref{b3:prop:integral-quotient-localization}，商包含 \(A/\mathfrak p\subseteq B/\mathfrak q\) 整。在此包含上对素理想 \(\mathfrak p'/\mathfrak p\) 应用\cref{b3:thm:lying-over}，得到 \(B/\mathfrak q\) 的素理想 \(\mathfrak r\)，其收缩为 \(\mathfrak p'/\mathfrak p\)。取 \(\mathfrak r\) 在 \(B\) 中的原像 \(\mathfrak q'\)，其在 \(A\) 中的收缩就是 \(\mathfrak p'/\mathfrak p\) 在 \(A\) 中的原像 \(\mathfrak p'\)，因而满足要求。
\end{proof}

从任意指定的底层原像出发，重复使用此定理，就能提升一个有限的上升素理想链。反向收缩一个严格链时，\cref{b3:thm:incomparability}保证严格性不会丢失。因此，整扩张两端的素理想链长可以相互比较。

\begin{proposition}\label{b3:prop:integral-spec-closed}
设 \(A\subseteq B\) 为交换含幺环的整扩张。它诱导的映射 \(\pi:\operatorname{Spec}B\to\operatorname{Spec}A\) 满射且为闭映射。更精确地，对理想 \(J\subseteq B\)，
\[
\pi\bigl(V_B(J)\bigr)=V_A(J\cap A).
\]
\end{proposition}
\begin{proof}
左侧的每个素理想显然包含 \(J\cap A\)。反向，对包含 \(J\cap A\) 的素理想，向整包含 \(A/(J\cap A)\subseteq B/J\) 应用\cref{b3:thm:lying-over}，便得到包含 \(J\) 的原像。取 \(J=0\) 给出满射性。
\end{proof}
